\documentclass[11pt]{amsart}
\usepackage{amssymb,mathtools}
\usepackage[margin=1.15in]{geometry}
\usepackage[hidelinks]{hyperref}
\newtheorem{theorem}{Theorem}[section]
\newtheorem{lemma}[theorem]{Lemma}
\newtheorem{proposition}[theorem]{Proposition}
\newtheorem{corollary}[theorem]{Corollary}
\theoremstyle{definition}\newtheorem{definition}[theorem]{Definition}
\theoremstyle{remark}\newtheorem{remark}[theorem]{Remark}
\newcommand{\OO}{\mathcal O}
\newcommand{\Spec}{\operatorname{Spec}}
\newcommand{\QC}{\operatorname{QC}}
\newcommand{\HH}{\mathrm H}
\newcommand{\CC}{\check{\mathrm H}}
\newcommand{\tr}[1]{\left\Vert#1\right\Vert}
\newcommand{\supp}{\operatorname{D}}
\newcommand{\code}[1]{{\small\texttt{#1}}}
\newcommand{\stamptime}{\the\numexpr(\time-30)/60\relax:\ifnum\numexpr\time-60*((\time-30)/60)\relax<10 0\fi\the\numexpr\time-60*((\time-30)/60)\relax}
\AddToHook{shipout/foreground}{\put(0.5\paperwidth,-\paperheight+0.55cm){\makebox[0pt]{\tiny\ttfamily unit/doc/\jobname.tex \ --- \ compiled \today, \stamptime}}}
\title[Constructive Serre affineness]{A constructive Serre criterion for affineness\\A detailed informal proof for locally ringed lattices}
\date{October 30, 2026}
\begin{document}\raggedbottom

\begin{abstract}
We give a point-free proof that a qcqs scheme whose quasi-coherent
modules have vanishing first cohomology is affine. For a finite affine
cover, add its support ideals and consider the fibre of their addition
map over $1$. Its translation module is a saturated submodule of a
finite product of support ideals, hence quasi-coherent. One application of the
torsor hypothesis gives a subordinate partition of unity. The resulting
global principal affine cover and the same unit equation give the final
affineness comparison. Canonical localisation maps are defined before
models; their equivalence descends through finite sheaf equalisers.
For the converse, the canonical maps of a quasi-coherent module are
equivalences on every open, by descent from the charts of an atlas; this
reconstructs the module on every affine open and reduces torsor
triviality to degree-one gluing. Vanishing of \v{C}ech $\HH^1$ on one
finite affine cover is equivalent to its vanishing on every such cover
and to affineness, and the groups of two such covers are canonically
isomorphic, through torsors. All of this is checked in the Unit
development \code{unit/sag/qcqs} (Section~\ref{sec:review}).
\end{abstract}
\maketitle

\section{Scope, conventions, and the theorem}
The geometric framework is the existing qcqs development in Unit.
A scheme is a small locally ringed distributive lattice
$X=(L,\OO)$ which merely has a finite affine open cover. The lattice
is bounded; its elements represent compact opens, its bottom is $0$,
and its top is $1$. All rings are unital commutative rings with set
carriers; the zero ring and the empty scheme are included.
Here and below ``scheme'' means qcqs scheme in this sense.

For $a\in L$ and $t\in\OO(a)$, write $D_a(t)$ for the greatest
$b\le a$ on which $t$ is invertible. Write $D(f)=D_1(f)$ for a
global section. Restriction is denoted $t|b$ when the source is clear.
Comparisons between different presentations of a ring or module use
their specified isomorphisms. An expression such as $\OO(D_A(h))=A_h$
below abbreviates its canonical comparison, not a judgemental equality.
All existential choices of presentations, denominators, covers, and
lifts are mere: the displayed constructions are first made for actual
witnesses and then eliminated only into propositions or uniquely
determined data. Finite choice for propositional truncation suffices.

\begin{theorem}[External constructive Serre criterion]\label{thm:serre}
Let $X$ be a small qcqs locally ringed lattice. Suppose that for every
small quasi-coherent $\OO_X$-module $M$, every $M$-torsor merely has a
global section. Then the canonical map
\[
 X\longrightarrow\Spec\OO_X(1)
\]
is an isomorphism. Conversely, if $X$ is affine, every torsor under a
small quasi-coherent module merely has a global section. Equivalently,
$X$ is affine if and only if $\CC^1(\mathcal U,M)=0$ for every small
quasi-coherent $M$, where $\mathcal U$ is one, or equivalently every,
finite affine cover (Section~\ref{sec:torsors}). For any $X$ with a
quasi-coherent $M$, the groups of two finite affine covers are
canonically isomorphic.
\end{theorem}

The direct implication uses the atlas-free torsor hypothesis as stated.
Its route is support ideals, a partition torsor, one
application of the hypothesis, and localisation descent for the
structure sheaf. Associated exactness, sheaf cokernels, general kernel
closure, affine reconstruction and the results of
Section~\ref{sec:torsors} are not prerequisites. Section~\ref{sec:modules}
provides the canonical localisation maps and a single criterion for
associated affine models. The converse
(Corollary~\ref{lem:affine-torsor}) and the $\HH^1$ form
(Corollary~\ref{cor:coverindependence}) are separate consequences.

The proof uses no higher cohomology, closed points, prime ideals,
generalized field axiom, synthetic scheme axiom, or resizing principle.
The original affine-replacement argument supplied in the conversation
was credited there to a discussion involving David and Felix. The
partition-of-unity proof below replaces its successive lifting steps
by one simultaneous lift. This note makes no further claim about priority or
the identities of the other participants.

\section{Geometric facts and finite affine refinements}
The following are interfaces already established in the geometric
development. They must be used by their names in the formal proof.
\begin{align}
 D_b(t|b)&=b\wedge D_a(t) &&(b\le a),\label{eq:restrict-support}\\
 D_a(1)&=a,&D_a(0)&=0,\label{eq:support01}\\
 D_a(st)&=D_a(s)\wedge D_a(t),&
 D_a(s+t)&\le D_a(s)\vee D_a(t).\label{eq:supportops}
\end{align}
Also $D_a(-t)=D_a(t)$, since $-1$ is a unit. A section $t$ is
invertible in $\OO(a)$ if and only if $D_a(t)=a$. For a finite cover
$a=\bigvee_i a_i$, the sheaf condition for units gives
\begin{equation}\label{eq:join-support}
 D_a(t)=\bigvee_i D_{a_i}(t|a_i).
\end{equation}
These formulas concern invertibility support; they do not assert that
it is the universal support of $\OO(a)$ for a general $a$.

On an affine open $w$, with $B=\OO(w)$, the canonical comparison
identifies the downset of $w$ with $\operatorname{Zar}(B)$. Every
$b\le w$ is therefore merely a finite join of $D_w(h_i)$.
Every such principal open is affine, and its ring of sections is
canonically $B_{h_i}$. In particular every compact open of a qcqs
scheme merely has a finite affine cover: intersect it with a finite
affine atlas and refine each intersection inside that atlas member.

\begin{lemma}[Common refinement]\label{lem:refine}
Two finite affine covers of a compact open merely admit a common
finite affine refinement, with specified inclusions into members of
each cover.
\end{lemma}
\begin{proof}
Intersect every pair of members. Each intersection is compact and is
an open of the first affine member, so express it as a finite join of
principal opens in that member. These are affine, lie in both original
members, and jointly cover the intersection. Collect the finitely many
finite families. This does not assert that the intersections themselves
are affine.
\end{proof}

One can also begin with the formulation in the question: a ringed
lattice with a finite affine cover, without including invertibility
support in its data. For a section $t\in\OO(a)$, intersect that cover
with $a$ and refine inside each affine chart by finitely many principal
opens. On each resulting affine piece the restricted section has an
invertibility open. Their finite join is an invertibility open on
$a$: a restriction is a unit exactly when its restrictions on this
finite cover are units, since inverses glue uniquely. Uniqueness of
the greatest such open makes the construction independent of all
refinements and presentations. The support laws are checked on the
affine pieces and then by the covering equation. Thus invertibility
support is uniquely recoverable from this affine-local condition.
For the subsequent proof we use the existing locally ringed record.

\section{Elementary module localisation}\label{sec:algebra}
The direct route uses fractions, finite sums and the finite-equaliser
localisation descent step. Cokernel exactness and degree-one gluing
belong to the broader algebraic and converse results recorded here.

For an $A$-module $N$, write $N_f$ for localisation at the powers of
$f$. Its elements are represented by $n/f^r$. Equality of two fractions
means merely that an additional power of $f$ annihilates their
cross-multiplied difference. In particular
\begin{equation}\label{eq:kill}
 n/1=0\text{ in }N_f\quad\Longleftrightarrow\quad
 \exists r\ge0\;f^r n=0.
\end{equation}
All maps of localised modules below are specified by their value on
$n/1$ and the inversion of $f$; their composition equations follow
by uniqueness, not by re-expanding fractions.

\begin{lemma}[Finite exactness]\label{lem:finite-exact}
Localisation of modules preserves finite products, kernels, and
cokernels, and hence finite equalisers and short exact sequences.
\end{lemma}
\begin{proof}
A finite tuple of fractions can be put over a common denominator,
giving the inverse to the product comparison. For a map $v:N\to P$,
if $v(n)/f^r=0$, equation~\eqref{eq:kill} gives $f^e v(n)=0$.
The fraction $n/f^r$ can then be written
$(f^e n)/f^{r+e}$ with numerator in $\ker v$. Injectivity follows
from the same equality criterion. For cokernels, lift a representative
of a quotient fraction to $N$ and use the fraction equality criterion
to identify its kernel with the image of the localised map. All these
representative choices are mere, and the induced maps are unique.
\end{proof}

\begin{lemma}[Powers of a unit-ideal family]\label{lem:powers}
If $(f_1,\ldots,f_m)=A$, then $(f_1^r,\ldots,f_m^r)=A$ for every
positive $r$.
\end{lemma}
\begin{proof}
The Zariski support has $D(f_i^r)=D(f_i)$, so the join of these opens
is top. The formal Nullstellensatz says that $1$ belongs to the ideal
of these powers. Equivalently, expand a sufficiently large power of
an equation $1=\sum_i a_i f_i$ and group terms by an index appearing
at least $r$ times. Thus coefficients $\lambda_i$ with
$\sum_i\lambda_i f_i^r=1$ merely exist.
\end{proof}

\begin{lemma}[Degree zero gluing]\label{lem:h0}
If $(f_1,\ldots,f_m)=A$, then
\[
 N\longrightarrow\prod_i N_{f_i}
 \rightrightarrows\prod_{i,j}N_{f_i f_j}
\]
is an equaliser.
\end{lemma}
\begin{proof}
For injectivity, a difference which vanishes in every $N_{f_i}$ is
killed by a power of each $f_i$. Choose one positive bound $r$ for
these finitely many powers. Lemma~\ref{lem:powers} then annihilates
the difference by $1$.

For existence, let $s_i\in N_{f_i}$ be compatible. Choose a common
positive $r$ and $t_i\in N$ with $t_i|f_i=f_i^r s_i$.
For every $i,j$, the difference
$t_i|f_j-f_i^r s_j$ vanishes after further localising at $f_i$.
Choose a common killing exponent $e$ for these finitely many
differences and replace $t_i$ by $f_i^e t_i$ and $r$ by $r+e$.
Now $t_i|f_j=f_i^r s_j$ in $N_{f_j}$ for all $i,j$.
Choose $\lambda_i$ with $\sum_i\lambda_i f_i^r=1$ and set
$t=\sum_i\lambda_i t_i$. Its image in $N_{f_j}$ is $s_j$.
Uniqueness has already been proved, so the mere representative choices
produce the uniquely determined amalgam.
\end{proof}

\begin{lemma}[Gluing over a basic open]\label{lem:h0relative}
Suppose $D_A(h_i)\le D_A(g)$ and
$D_A(g)\le\bigvee_i D_A(h_i)$. Then
\[
 N_g\longrightarrow\prod_i N_{h_i}
 \rightrightarrows\prod_{i,j}N_{h_i h_j}
\]
is an equaliser, with the canonical localisation restrictions.
\end{lemma}
\begin{proof}
The denominator and error clearing in Lemma~\ref{lem:h0} take place
in the original module and its localisations, before the partition of
one is used. They therefore give one exponent $r$ and numerators
$t_i\in N$ satisfying $t_i|h_j=h_i^r s_j$ for compatible local data
$s_j\in N_{h_j}$. Since $D_A(h_i)\le D_A(h_i^r)$, the formal
Nullstellensatz gives, merely,
\[
 g^e=\sum_i\lambda_i h_i^r.
\]
Put $t=\sum_i\lambda_i t_i$ and take $s=t/g^e\in N_g$.
In $N_{h_j}$ we have $t|h_j=g^e s_j$, and $g$ acts invertibly
there because $D_A(h_j)\le D_A(g)$. Hence $s|h_j=s_j$.

For uniqueness, represent a difference by $x/g^q$. Vanishing on
every $N_{h_i}$ implies that powers of the $h_i$ kill $x$.
The formal Nullstellensatz expresses a power of $g$ as a linear
combination of those killing powers. Thus a power of $g$ kills $x$,
so $x/g^q=0$. All choices are eliminated into equality or the
proposition that the amalgam exists uniquely. The argument includes
an empty covering family and the zero ring.
\end{proof}

\subsection*{Degree-one gluing for the converse}
This final algebraic lemma is used for affine torsor triviality,
not for the direct implication of the criterion.

\begin{lemma}[Degree one gluing]\label{lem:h1affine}
If $\bigvee_iD_A(f_i)=1$, every family
$c_{ij}\in N_{f_i f_j}$ satisfying
\[
 c_{ij}|f_k+c_{jk}|f_i=c_{ik}|f_j
 \quad\text{in }N_{f_i f_j f_k}
\]
merely admits $b_i\in N_{f_i}$ with $b_j-b_i=c_{ij}$ on every
overlap.
\end{lemma}
\begin{proof}
Choose a common positive $r$ and elements $a_{ki}\in N_{f_i}$
whose images in $N_{f_k f_i}$ are $f_k^r c_{ki}$. This is possible
by clearing the $f_k$ denominators in the finitely many $c_{ki}$.
In $N_{f_i f_j}$ put
\[
 e_{kij}=a_{kj}|f_i-a_{ki}|f_j-f_k^r c_{ij}.
\]
After localising at $f_k$ this is zero by the cocycle equation.
Choose one exponent $e$ killing every $e_{kij}$ by the appropriate
$f_k^e$. Replace $a_{ki}$ by $f_k^e a_{ki}$ and $r$ by $r+e$.
We now have the actual equations
\[
 a_{kj}|f_i-a_{ki}|f_j=f_k^r c_{ij}
 \quad\text{in }N_{f_i f_j}.
\]
The lattice covering equation, the equality $D_A(f_k^r)=D_A(f_k)$
and the formal Nullstellensatz give, merely, $\lambda_k$ with
$\sum_k\lambda_k f_k^r=1$. Set
$b_i=\sum_k\lambda_k a_{ki}$. Subtracting the two expressions on
an overlap gives $b_j-b_i=c_{ij}$. All choices remain inside the
mere assertion that a coboundary exists.
\end{proof}

For an empty covering family, $A$ is the zero ring and every unital
$A$-module is zero. Both lemmas then assert contractibility of the
appropriate empty diagram, and hold as well.

\section{Module sheaves and quasi-coherence}\label{sec:modules}
\begin{definition}
A module sheaf $M$ on $X$ assigns an $\OO(a)$-module $M(a)$ to
each $a\in L$, with additive restriction maps which respect scalar
restriction. It satisfies the finite-cover sheaf condition on the
lattice. A morphism is a natural family of linear maps. It is locally
surjective if for every $y\in F(a)$ there merely exist a finite cover
$a=\bigvee_i a_i$ and $x_i\in E(a_i)$ mapping to $y|a_i$.
\end{definition}

For an $A$-module $N$, define $\widetilde N$ on $\Spec A$ directly:
on an open $w$, a section is a compatible family of elements of $N_h$
for all $D_A(h)\le w$, with compatibility under basic-open restriction.
The construction is the module analogue of the existing construction
of the affine structure sheaf. Evaluation on the component $h$ gives
\begin{equation}\label{eq:module-basic}
 \widetilde N(D_A(h))\cong N_h.
\end{equation}
Its inverse restricts an element of $N_h$ to every basic open below
$D_A(h)$. The inverse laws are componentwise compatibility and
uniqueness of localisation maps. For a general finite cover, refine
inside basic opens and use Lemma~\ref{lem:h0relative} for the resulting
cover of each basic open. This gives sheafhood. Ring operations, module operations,
and scalar restriction are componentwise. A module map induces a
sheaf morphism by its localisations. These observations also prove
\[
 \widetilde N|D_A(h)\cong\widetilde{N_h}
 \quad\text{over }A_h.
\]

\begin{definition}\label{def:qc}
A module sheaf $M$ is quasi-coherent if there merely exist a finite
affine cover $w_i$ of $X$ and modules $N_i$ over $\OO(w_i)$ with
$M|w_i\cong\widetilde{N_i}$. No chosen atlas or chosen $N_i$ is
part of the property.
\end{definition}

Restriction of a quasi-coherent module to a compact open is
quasi-coherent: refine its intersections with the given $w_i$ by
principal opens inside $w_i$ and use the last displayed comparison.
The same argument proves locality: if a finite compact-open cover has
quasi-coherent restrictions, collecting their finite affine witnesses
gives witnesses for Definition~\ref{def:qc}. Both conclusions are
propositions, so the atlas choices may be eliminated. Neither fact is
used in the proof of Theorem~\ref{thm:serre}.

\subsection*{Canonical maps first; models only to prove equivalence}
For any module sheaf $M$, any $a\in L$ and $g\in\OO(a)$, set
$a_g=D_a(g)$. Restriction gives an $\OO(a)$-linear map
$M(a)\to M(a_g)$, with target scalars restricted along
$\OO(a)\to\OO(a_g)$. Since $g|a_g$ is a unit, its action on
$M(a_g)$ is invertible. The module-localisation universal property
therefore gives a specified map
\begin{equation}\label{eq:canonical-module-localisation}
 \lambda_{a,g}:M(a)_g\longrightarrow M(a_g),\qquad
 \lambda_{a,g}(x/g^r)=(g|a_g)^{-r}(x|a_g).
\end{equation}
No quasi-coherence or local model is used in this definition.
For $b\le a$, restriction and localisation give a commuting square
\[
\begin{array}{ccc}
 M(a)_g&\xrightarrow{\lambda_{a,g}}&M(D_a(g))\\
 \downarrow&&\downarrow\\
 M(b)_{g|b}&\xrightarrow{\lambda_{b,g|b}}&M(D_b(g|b)).
\end{array}
\]
Indeed, after restriction of scalars to $\OO(a)$ both composites are
linear extensions of the same restriction on original sections; the
action of $g$ is invertible in the target. Uniqueness of localisation
extensions proves their equality. The same argument gives compatibility
when one basic open $D(h)$ is contained in another $D(g)$: the action
of $g$ is already invertible on $D(h)$. Thus every comparison diagram below
is built from canonical maps which already commute. Local models are
used only to prove the components are equivalences.

\begin{lemma}[The $\lambda$-criterion]\label{lem:lambda-criterion}
Let $w$ be affine and put $B=\OO(w)$. The following are equivalent:
$M|w$ is associated to a $B$-module; and for every $s\in B$ the
canonical map
\[
 \lambda_{w,s}:M(w)_s\longrightarrow M(D_w(s))
\]
is an isomorphism. Under the second condition the maps extend uniquely
to a natural linear isomorphism
$\widetilde{M(w)}\cong M|w$, along the canonical affine chart. Under
the first, more generally $\lambda_{v,s|v}$ is an isomorphism for every
principal open $v=D_w(h)$ and every $s\in B$.
\end{lemma}
\begin{proof}
Both directions use the element form of a localisation equivalence.
For a $B$-linear $u:P\to Q$ such that $f$ acts invertibly on $Q$, the
extension $P_f\to Q$ is bijective if and only if $u$ \emph{clears}
(every $y\in Q$ merely has $f^ry=u(x)$) and \emph{kills} ($u(x)=0$
implies merely $f^rx=0$). Both conditions are elementwise, so they pass
along a commuting square whose sides are bijective $B$-linear maps, and
they are unchanged when an $\OO(v)$-module is regarded as a $B$-module
along $B\to\OO(v)$, with $s$ in place of its image.

Let $e$ be a model, $M|w\cong\widetilde N$ along the chart $\pi$, let
$v=D_w(h)=\pi(D_B(h))$ and put $t=sh$. Then
$D_v(s|v)=v\wedge D_w(s)=D_w(sh)=\pi(D_B(t))$. Evaluating $e$ on these
basic opens gives bijections $i:M(v)\cong N_h$ and
$j:M(D_v(s|v))\cong N_t$. They are $B$-linear: the chart sends the
section $b/1$ to $b|$, $e$ is linear over the structure sheaf, and
evaluation is semilinear over ring evaluation. Naturality of $e$ and of
evaluation turns restriction $\rho:M(v)\to M(D_v(s|v))$ into the
restriction $N_h\to N_{sh}$, that is $j\rho=ri$. The map $r$ clears
and kills at $s$: the restriction of $x/h^p$ is $s^px/(sh)^p$, and a
fraction vanishing in $N_{sh}$ is killed by a power of $s$ in $N_h$.
Hence $\rho$ clears and kills, and $\lambda_{v,s|v}$ is an isomorphism.
The case $h=1$, $v=w$, is the first direction. Affineness of $w$ is
not used.

Conversely, evaluate $\widetilde{M(w)}$ on the principal basis and
use $\lambda_{w,s}$ as its comparison with $M$. The squares proved
above give naturality, including arbitrary inclusions between basic
opens. Regard the maps as linear maps over the structure sheaf using
the canonical basic-ring comparisons. The basis-map extension theorem
constructs their unique global morphism. The local-equivalence theorem,
using finite basic covers, makes that morphism an equivalence.
No section-localisation theorem or reconstruction result is assumed.
\end{proof}

Consequently Definition~\ref{def:qc} is equivalently the mere existence
of a finite affine cover on whose members all canonical $\lambda$ maps
are isomorphisms. We retain the associated-model definition, but use
this criterion to supply its models. Finite choices of models or covers
remain under truncation until the quasi-coherence property is obtained.

\begin{lemma}[Finite sums]\label{lem:finite-sums}
Finite direct sums of module sheaves are formed sectionwise and are
sheaves. On an affine $w$, if $M_i|w\cong\widetilde{N_i}$, then
\[
 (\bigoplus_iM_i)|w\cong\widetilde{\bigoplus_iN_i}.
\]
The empty sum is the zero sheaf.
\end{lemma}
\begin{proof}
For a finite family, sums are products. Compatible tuples glue
coordinatewise, including the empty tuple. Localisation commutes with
finite products by choosing common denominator bounds; the canonical
maps for the sum are the coordinatewise canonical maps. Apply the
$\lambda$-criterion, or equivalently the same basic comparisons, to
obtain the displayed associated model. Its coefficient module is the
finite sum of the $N_i$, using their top evaluations.
\end{proof}

\begin{theorem}[Descent of canonical localisation]
\label{thm:sectionloc}\label{thm:localisation-descent}
Let $u_i$ be a finite cover of $X$, $M$ a module sheaf and
$g\in A=\OO(1)$. Suppose the canonical localisation maps on every
member $u_i$ and every overlap $u_i\wedge u_j$, at the respective
restrictions of $g$, are isomorphisms. Then
\[
 \lambda_{1,g}:M(1)_g\longrightarrow M(D(g))
\]
is an isomorphism of $A_g$-modules.
Neither affineness nor associated models are assumptions of this
pure descent statement.
\end{theorem}
\begin{proof}
The sheaf condition expresses $M(1)$ as the equaliser
\[
 M(1)\longrightarrow\prod_iM(u_i)
 \rightrightarrows\prod_{i,j}M(u_i\wedge u_j).
\]
Regard every term as an $A$-module through restriction. Localise at
$g$ and compare with the sheaf equaliser for the cover of $D(g)$ by
$u_i\wedge D(g)$. All maps are the canonical maps defined before any
models; their squares commute by the localisation universal property.
Localisation over $A$ agrees with localisation at the image of $g$
in each section ring, after explicit restriction of scalars. The target
diagram is also regarded as an $A$-module diagram, by restriction from
$\OO(D(g))$.

The finite-equaliser localisation descent lemma, with the assumed
member and overlap equivalences, makes the middle canonical map an
isomorphism. Its algebraic input is preservation of finite products
and equalisers under localisation, the relevant part of
Lemma~\ref{lem:finite-exact}. No additional kernel or associated-exactness
calculation is required here. The empty-cover sheaf condition gives
zero modules on the empty space and includes that case.
\end{proof}

\begin{corollary}[Canonical localisation on an affine]
\label{thm:principal-module-localisation}\label{cor:affine-basic-sections}
Let $X$ be affine, $A=\OO(1)$ and $M$ quasi-coherent. Then
$\lambda_{1,g}:M(1)_g\to M(D(g))$ is an isomorphism for every
$g\in A$. Consequently $M\cong\widetilde{M(1)}$ along the canonical
chart (affine reconstruction).
\end{corollary}
\begin{proof}
The claim is a proposition; merely take quasi-coherent witnesses
$w_j$ with models. The canonical lattice map
$\operatorname{Zar}(A)\to L$ is bijective, so the preimages of the
$w_j$ cover $\operatorname{Zar}(A)$ and are merely refined by finitely
many $D_A(f_k)$, with $D(f_k)\le w_{j(k)}$ and $\bigvee_kD(f_k)=1$.
The members $D(f_k)$ and the overlaps $D(f_k)\wedge D(f_l)=D(f_kf_l)$
are principal opens $D(f)$ contained in the chart $w=w_{j(k)}$, and
such an open is the principal open $D_w(f|w)$ of the chart. By the
$\lambda$-criterion, in its general form, $\lambda$ at $D(f)$ at the
restriction of $g$ is an isomorphism, since $g|D(f)=(g|w)|D(f)$.
Theorem~\ref{thm:localisation-descent} gives $\lambda_{1,g}$. The top
open is affine, so the $\lambda$-criterion at $w=1$ gives the model.
No restriction of models or of quasi-coherence to smaller opens is
used.
\end{proof}

\begin{corollary}[Canonical localisation on every open]\label{cor:lambda-everywhere}
If $M$ is quasi-coherent, then $\lambda_{w,g}$ is an isomorphism for
every open $w$ and every $g\in\OO(w)$. Hence on every affine open $w$,
$M|w\cong\widetilde{M(w)}$ along the chart.
\end{corollary}
\begin{proof}
Say $M$ is good at $w$ when every $\lambda_{w,g}$ is an isomorphism.
(1) Let $w$ be a chart with a model and $v=D_w(h)$. Every $c\in\OO(v)$
is $e\,(x|v)$ with $e$ a unit: the ring map $B_h\to\OO(v)$ is
bijective, so $c$ is the image of some $x/h^p$, and $h^pc=x|v$ with
$h|v$ a unit. The $\lambda$-criterion gives $\lambda_{v,x|v}$; a unit
factor changes neither $D_v$ nor clearing and killing, so
$\lambda_{v,c}$ is an isomorphism. (2) If $u\le w$ for such a chart,
the preimage of $u$ in $\operatorname{Zar}(B)$ is merely
$D_B(h_1)\vee\dots\vee D_B(h_n)$; so $u$ is covered by the chart opens
$D_w(h_i)$, whose meets $D_w(h_ih_k)$ are chart opens, and pure descent
on $X|u$ shows $M$ good at $u$. (3) For any $w$ and an atlas $w_j$, the
$w\wedge w_j$ cover $w$, and they and their meets lie below the charts:
descent on $X|w$ and (2). The model on an affine $w$ is the
$\lambda$-criterion. Each descent is on a slice, and $\lambda$ of an
open of a slice is $\lambda$ of that open in $X$.
\end{proof}

For the direct implication, apply pure descent to $\OO$ using only
affine chart ring facts. The corollary is used in the converse, where
the model it supplies reduces affine torsor triviality to associated
modules. No compact-open localisation lemma or reconstruction sheaf map
is needed for either application.

\subsection*{Optional module theory}
The following results are consequences, not prerequisites for the
criterion or for affine torsor triviality.
\begin{lemma}[Exactness for associated modules]\label{lem:associated-exact}
For an $A$-linear map $u:N\to P$, there are canonical comparisons
\[
 \widetilde{\ker u}\cong\ker(\widetilde u),\qquad
 \widetilde{\operatorname{coker}u}\cong
 \operatorname{coker}_{\mathrm{sheaves}}(\widetilde u).
\]
The map $\widetilde P\to\widetilde{\operatorname{coker}u}$ is locally
surjective. In particular, for an ideal $I\subseteq A$, the kernel of
$\widetilde A\to\widetilde{A/I}$ is $\widetilde I$.
\end{lemma}
\begin{proof}
The sectionwise kernel of a module-sheaf map is a sheaf: glue in the
source and check the kernel equation by separation in the target.
On $D_A(h)$, Lemma~\ref{lem:finite-exact} identifies
$(\ker u)_h$ with $\ker(u_h)$. Evaluation
\eqref{eq:module-basic} gives the comparison with the sheaf kernel
there. These maps commute with restriction, since they commute with
the kernel inclusions into $N_h$. Extend over finite basic covers by
sheaf gluing and separation to obtain the first comparison.

For the quotient, finite exactness identifies
$(\operatorname{coker}u)_h$ with $\operatorname{coker}(u_h)$.
An element of this quotient merely lifts to $P_h$. Thus the map to
the associated quotient is locally surjective, by a finite basic
refinement of any compact open. A sheaf morphism from $\widetilde P$
which kills $\widetilde u$ factors uniquely on each basic open through
this module quotient; uniqueness makes the factors natural, and the
sheaf condition extends them to all opens. This proves the sheaf
cokernel property. For the last assertion apply the kernel comparison
to $A\to A/I$, whose module kernel is $I$.
\end{proof}

This lemma concerns maps already supplied as associated maps. It
requires neither affine reconstruction nor localisation of sections
on a general compact open.


\begin{corollary}\label{cor:affine-reconstruction}
On an affine $X=\Spec A$, every quasi-coherent module $M$ is
canonically $\widetilde{M(1)}$. A morphism between such sheaves is
determined on basic opens by the localisations of its global map.
\end{corollary}
\begin{proof}
Corollary~\ref{cor:affine-basic-sections} gives the comparison on
every basic open, with its restriction equations. Extend it to all opens using
their finite basic covers and the sheaf condition. For a morphism,
linearity and naturality force the value on $n/f^r$ to be the image
of $n$ divided by $f^r$; these fractions exhaust the localised module.
\end{proof}

\begin{lemma}[Kernels and quotient sheaves]\label{lem:kernels}
The kernel of a morphism between quasi-coherent modules is
quasi-coherent. For a module map $N\to P$ on an affine, the sheaf
$\widetilde{\operatorname{coker}(N\to P)}$ is its sheaf cokernel.
The map to this cokernel is locally surjective.
\end{lemma}
\begin{proof}
The cokernel assertions are Lemma~\ref{lem:associated-exact}.
For the kernel, take finite affine witnesses for the two quasi-coherent
sheaves and a common finite affine refinement. On that refinement,
Corollary~\ref{cor:affine-reconstruction} identifies both sheaves and
the morphism with their associated global modules and map.
Lemma~\ref{lem:associated-exact} then identifies the sheaf kernel
with the associated module kernel on each member. These models give
quasi-coherence. This general closure result uses reconstruction;
the partition kernel in the direct proof does not.
\end{proof}

\section{Torsors and first cohomology}\label{sec:torsors}
The torsor definition below is used in the hypothesis and the lifting
lemma. The results of this section establish the converse and the
cohomological reformulation; the direct implication uses no affine
torsor-triviality or classification theorem.

An $M$-torsor $T$ is a sheaf of sets with an action of the additive
sheaf $M$. For every $a$, whenever $t,t'\in T(a)$ there is a unique
$m\in M(a)$ with $t'=t+m$; and every $a$ merely has a finite
cover on which $T$ has sections. Equivalently, the action-difference
map $M\times T\to T\times T$ is an isomorphism and $T$ is locally
inhabited. A global section identifies $T$ with $M$ by translation.
An equivariant morphism of torsors is automatically an isomorphism,
as is checked after a local trivialisation and then by sheafhood.

\begin{lemma}[Torsors under an associated module on an affine]
\label{lem:associated-affine-torsor}
For an $A$-module $N$, every torsor under $\widetilde N$ on
$\Spec A$ merely has a global section.
\end{lemma}
\begin{proof}
Local inhabitation at the top gives a finite compact-open cover with
sections. Refine it by finitely many principal opens $D_A(f_i)$ and
restrict the sections to obtain $t_i\in T(D_A(f_i))$, with
$\bigvee_iD_A(f_i)=1$.

On each overlap write $t_i=t_j+c_{ij}$. Evaluation at $f_if_j$
gives $c_{ij}\in N_{f_if_j}$, and uniqueness of differences gives
the cocycle equation on triple overlaps. Apply
Lemma~\ref{lem:h1affine} directly with the lattice covering equation
to obtain, merely, $b_i\in N_{f_i}$ with $b_j-b_i=c_{ij}$.
Extend $b_i$ to the corresponding associated section on $D_A(f_i)$.
The corrected sections $t_i+b_i$ agree on overlaps and therefore glue.
All choices remain inside the mere existence assertion. If the cover
is empty, the same sheaf gluing step supplies the unique section.
No separate conversion of the cover to a unit-ideal witness is needed
in this torsor argument; that algebra is internal to degree-one gluing.
\end{proof}

\begin{lemma}[Torsors along a chart]\label{lem:torsor-chart}
Let $w$ be affine, $B=\OO(w)$, and $e$ a model of $M|w$: an
isomorphism from the direct image of $M|w$ along the chart
$\pi:\operatorname{Zar}(B)\to L$ to $\widetilde N$. For an $M$-torsor
$T$ put $T'(V)=T(\pi V)$, with $t+x=t+e^{-1}(x)$. Then $T'$ is a
$\widetilde N$-torsor on $\Spec B$, and a global section of $T'$ is a
section of $T$ over $w$.
\end{lemma}
\begin{proof}
The chart preserves finite joins and meets, so a cover of $V$ gives one
of $\pi V$, compatibility transfers along the meet equation, and the
sheaf condition is that of $T$. The action laws are those of $T$, since
$e^{-1}$ is additive and natural. The difference of $t,u$ is
$e(u-t)$, uniquely, because $e$ is bijective. For local sections, cover
$\pi V$ by opens with sections; the chart is bijective since $w$ is
affine, so these opens are images of opens of $\Spec B$, and since $\pi$
preserves meets and is injective, the order of $\Spec B$ is reflected:
the preimages lie below $V$ and cover it. Finally $\pi(1)=w$.
\end{proof}

\begin{corollary}[Affine torsor triviality]\label{lem:affine-torsor}
Every torsor under a quasi-coherent module on an affine scheme
merely has a global section. More generally, on any $X$, every torsor
under a quasi-coherent module merely has a section over every affine
open.
\end{corollary}
\begin{proof}
By Corollary~\ref{cor:affine-basic-sections}, $M$ has a model
$M\cong\widetilde{M(1)}$ along the chart of the top open. Move the
torsor along it (Lemma~\ref{lem:torsor-chart}) and apply
Lemma~\ref{lem:associated-affine-torsor}. The empty case is included.
For an affine open $w$, use instead the model at $w$ of
Corollary~\ref{cor:lambda-everywhere}; Lemma~\ref{lem:torsor-chart}
then gives a section over $w$.
\end{proof}

Fix a finite cover $\mathcal U=(u_i)$ of $X$. Define
$Z^1(\mathcal U,M)$ to be families $c_{ij}\in M(u_i\wedge u_j)$
with $c_{ij}+c_{jk}=c_{ik}$ after restriction to triple overlaps.
These equations imply $c_{ii}=0$ and $c_{ji}=-c_{ij}$.
For $b_i\in M(u_i)$ put
$(\delta b)_{ij}=b_j|u_i-b_i|u_j$.
Define $c\sim c'$ when there merely exist $b_i$ with
$c'=c+\delta b$. This is an equivalence relation compatible with
addition, so the set quotient
\[
 \CC^1(\mathcal U,M)=Z^1(\mathcal U,M)/\!\sim
\]
is an abelian group. The cocycle type and this quotient are small.

\begin{proposition}[Cocycles and torsors]\label{prop:classification}
The group $\CC^1(\mathcal U,M)$ classifies, up to equivariant
isomorphism, the $M$-torsors which merely have sections on every $u_i$.
For such torsors, zero corresponds exactly to those which merely have a
global section.
\end{proposition}
\begin{proof}
For $c\in Z^1(\mathcal U,M)$, define $T_c(a)$ to be families
\[
 x_i\in M(a\wedge u_i),\qquad
 x_j-x_i=c_{ij}|(a\wedge u_i\wedge u_j).
\]
Restriction and the $M(a)$-action are componentwise. This is a sheaf:
glue each coordinate and use separation for the displayed equations.
If $a\le u_i$, the coordinate $x_i$ determines all others by
$x_j=x_i|u_j+c_{ij}|(a\wedge u_j)$. The cocycle law proves
compatibility. Thus $T_c|u_i$ is trivial. For any two points,
their coordinate differences agree on overlaps and glue to the unique
element of $M(a)$ translating one to the other. Hence $T_c$ is a
torsor. The local point with $i$-coordinate zero will be denoted
$t_i$; its differences satisfy $t_i=t_j+c_{ij}$.

Replacing $c$ by $c+\delta b$ gives the equivariant isomorphism
$(x_i)\mapsto(x_i+b_i)$. Conversely any equivariant isomorphism
between two such torsors is translation by a unique $b_i$ in their
local coordinates; its compatibility equations give exactly the
coboundary relation. Consequently the quotient identifies precisely
the isomorphism classes, not any extra pairs.

For a torsor $T$ with sections on every $u_i$, finite choice gives a
mere family $t_i$. Their differences $t_i=t_j+c_{ij}$
give a cocycle. A point $t\in T(a)$ has coordinates $x_i$ defined
by $t|u_i=t_i+x_i$; then $x_j-x_i=c_{ij}$.
This identifies $T$ with $T_c$. Changing the local sections changes
the cocycle by a coboundary, so its quotient class is uniquely
determined and can be extracted from the mere family of sections.
If $T$ has a global section, its coordinates show the class is zero;
if the class is zero, the coboundary equation adjusts the $t_i$ to
a compatible family and glues a global section, merely.
\end{proof}

\begin{corollary}[The \texorpdfstring{$\HH^1$}{H1} form]\label{cor:coverindependence}
Let $\mathcal U=(u_i)$ be a finite cover of $X$ and $M$ a module sheaf.
\begin{enumerate}
\item If every $M$-torsor merely has a global section, then
 $\CC^1(\mathcal U,M)=0$.
\item If $\CC^1(\mathcal U,M)=0$, every $M$-torsor with sections on the
 members $u_i$ merely has a global section.
\item If $\mathcal U$ is affine and $\CC^1(\mathcal U,K)=0$ for every
 quasi-coherent $K$, then $X$ is affine.
\end{enumerate}
Hence, for a qcqs $X$, the following are equivalent: $X$ is affine;
$\CC^1(\mathcal U,M)=0$ for every finite affine cover $\mathcal U$ and
every quasi-coherent $M$; this holds for one merely given finite affine
cover.
\end{corollary}
\begin{proof}
(1) Every class is that of a cocycle $c$; the class of $T_c$ is $c$,
and a torsor with a global section has class zero.
(2) The class of the torsor is zero; apply
Proposition~\ref{prop:classification}.
(3) The partition torsor of Lemma~\ref{lem:lifts} for the cover
$\mathcal U$ has points on the members: the basis vector of $1$ on
$u_j$. Its translation module is quasi-coherent
(Proposition~\ref{prop:principalcover}), so by (2) the torsor merely has
a global section, which is a partition of unity subordinate to
$\mathcal U$; conclude as in Sections~\ref{sec:partition}
and~\ref{sec:finish}. For the equivalence, an affine $X$ satisfies the
torsor hypothesis by Corollary~\ref{lem:affine-torsor}, so (1) applies
to every cover; a qcqs $X$ merely has a finite affine cover; and (3)
concludes.
\end{proof}

Thus the theorem is stated by vanishing, and its $\HH^1$ form needs no
comparison of the groups of two covers. That comparison nevertheless
holds, through torsors.

\begin{lemma}[Refinement]\label{lem:refinement}
Let $\mathcal W=(w_a)$ refine $\mathcal U$, with $w_a\le u_{\tau a}$.
Restriction $\rho(c)_{ab}=c_{\tau a\tau b}|(w_a\wedge w_b)$ induces an
additive map $\CC^1(\mathcal U,M)\to\CC^1(\mathcal W,M)$ which sends the
$\mathcal U$-class of every torsor to its $\mathcal W$-class. If every
$M$-torsor has sections on the members of $\mathcal U$, it is bijective.
\end{lemma}
\begin{proof}
The restricted cochain is a cocycle: restrict the cocycle law of
$\mathcal U$ to the triple overlap. For sections $t_i$ on $\mathcal U$,
the restricted sections $t_{\tau a}|w_a$ have cocycle $\rho$ of that of
the $t_i$, since differences commute with restriction. In particular
$\rho(c)$ is the cocycle of the restricted canonical points of $T_c$.
If $c\sim d$, then $T_c\cong T_d$ (Proposition~\ref{prop:classification}
on $\mathcal U$), so their $\mathcal W$-classes agree; this makes the map
well defined, and additivity is that of restriction. Under the
hypothesis, send $[d]$ to the $\mathcal U$-class of $T_d$. Its image is
the $\mathcal W$-class of $T_d$, that is $[d]$; and $T_{\rho(c)}$, $T_c$
have the same $\mathcal W$-class, so are isomorphic and have the same
$\mathcal U$-class $[c]$.
\end{proof}

\begin{corollary}[Independence of the affine cover]\label{cor:groupcomparison}
For a quasi-coherent $M$ and finite affine covers $\mathcal U$,
$\mathcal V$, there is a canonical isomorphism
$\CC^1(\mathcal U,M)\cong\CC^1(\mathcal V,M)$ of groups. It sends the
$\mathcal U$-class of every torsor to its $\mathcal V$-class, in
particular $[c]$ to the class of $T_c$, and its inverse is the
comparison from $\mathcal V$ to $\mathcal U$.
\end{corollary}
\begin{proof}
Let $\mathcal Y$ be the concatenation of the two families; $\mathcal U$
and $\mathcal V$ refine it with equal members. Every torsor has
sections on the affine members of $\mathcal Y$
(Corollary~\ref{lem:affine-torsor}), so Lemma~\ref{lem:refinement}
gives additive bijections $\CC^1(\mathcal Y)\to\CC^1(\mathcal U)$ and
$\CC^1(\mathcal Y)\to\CC^1(\mathcal V)$, both sending the class of a
torsor to its class. Compose the second with the inverse of the first.
No contracted sum of torsors and no common refinement by intersections
is needed, and the restriction of quasi-coherence to affine opens is
replaced by Corollary~\ref{cor:lambda-everywhere}.
\end{proof}

The atlas-free proposition
\[
 \operatorname{Acyclic}_1(X,M):\quad
 \text{every $M$-torsor merely has a global section}
\]
is the hypothesis of the formal proof. This avoids extracting a chosen
atlas from the qcqs property in order to state the hypothesis.
The classification uses a small Cech quotient to represent torsor
classes; it does not postulate a small type of all module sheaves or
all torsors. Quantifying over those structures may live one universe
higher, as in the existing geometric category construction.

No affineness is used in Proposition~\ref{prop:classification}, and
only the members of $\mathcal U$ in Corollary~\ref{cor:coverindependence}(3).
The pair and triple intersections carry section groups and equations;
their affineness and their higher acyclicity are not used.

\section{The fibre of an addition map}\label{sec:lifting}
\begin{lemma}[Partition torsor]\label{lem:lifts}
Let $u_1,\ldots,u_m$ be a finite cover and let $J_i$ be sheaves of
ideals of $\OO$ such that $J_i|u_i=\OO|u_i$. Define sectionwise
\[
 E=\bigoplus_iJ_i,\qquad
 K(a)=\{(t_i)\in E(a)\mid\sum_i t_i=0\},\qquad
 T(a)=\{(t_i)\in E(a)\mid\sum_i t_i=1\}.
\]
Then $K$ is a module sheaf and $T$ is a $K$-torsor. If $K$ is
quasi-coherent, the hypothesis of Theorem~\ref{thm:serre} gives,
merely, $t_i\in J_i(1)$ with $\sum_i t_i=1$.
\end{lemma}
\begin{proof}
Finite tuples glue coordinatewise. Each displayed equation can be
checked locally by separation in $\OO$, so both fibres are sheaves.
The homogeneous equation defining $K$ is preserved by the sectionwise
module operations, making $K$ a module sheaf.
Addition of a zero-sum tuple acts on $T$. For $t,t'\in T(a)$,
the unique translating element is $t'-t\in K(a)$.
On $u_i$, the tuple with $i$th coordinate $1$ and all others $0$
is a section of $T$. Restrictions of these sections to $a\wedge u_i$
prove local inhabitation on every $a$. Thus $T$ is a $K$-torsor,
and the atlas-free hypothesis applies once $K$ is quasi-coherent.
For $m=0$ the cover implies $1_L=0_L$; the empty-cover sheaf condition
gives $1=0$ in the section ring and the unique empty tuple.
\end{proof}

Only these local sections of this particular fibre are needed.
No theorem about local surjectivity on arbitrary targets is required.
The existing general sheaf-of-lifts construction may be reused for
this special case, but it is not a further mathematical prerequisite.

\section{The support ideal}\label{sec:supportideal}
Fix $U\in L$. Define
\begin{equation}\label{eq:J}
 J_U(a)=\{t\in\OO(a)\mid D_a(t)\le U\}.
\end{equation}
The section $t$ is the underlying element; membership is a proposition.

\begin{lemma}[Ideal sheaf]\label{lem:J-sheaf}
$J_U$ is a sheaf of ideals in $\OO$ and $J_U|U=\OO|U$ canonically.
\end{lemma}
\begin{proof}
Zero belongs by~\eqref{eq:support01}. If $s,t$ belong, then
$D_a(s+t)\le D_a(s)\vee D_a(t)\le U$; negatives belong because
$D_a(-t)=D_a(t)$. For $r\in\OO(a)$,
$D_a(rt)=D_a(r)\wedge D_a(t)\le U$. Thus $J_U(a)$ is an ideal.
Equation~\eqref{eq:restrict-support} proves restriction stability.
For a compatible family of ideal sections, first glue in $\OO$;
equation~\eqref{eq:join-support} proves that the glued section
belongs to $J_U$. Uniqueness follows from uniqueness in $\OO$,
because ideal membership is a proposition. If $a\le U$, every
$D_a(t)$ is below $a\le U$, so the inclusion $J_U(a)\to\OO(a)$
is bijective, compatibly with all operations and restrictions.
\end{proof}

\begin{lemma}[Saturation]\label{lem:saturation}
Let $S\subseteq E$ be a submodule sheaf, $a\in L$ and $g\in\OO(a)$,
and suppose $\lambda^E_{a,g}$ is an isomorphism. Suppose $S$ is
\emph{saturated at $g$}: every $x\in E(a)$ with $x|D_a(g)\in S(D_a(g))$
merely has some $e$ with $g^ex\in S(a)$. Then $\lambda^S_{a,g}$ is an
isomorphism.
\end{lemma}
\begin{proof}
Write $\iota$ for the inclusion. Its localisation $\iota_g$ is
injective: if $\iota(x)/g^r=0$ then $g^n x=0$ in $E(a)$, hence in
$S(a)$. By definition of the canonical maps,
$\lambda^E\circ\iota_g=\iota\circ\lambda^S$, so $\lambda^S$ is
injective. For $y\in S(D_a(g))$ write $\iota y=\lambda^E(x/g^r)$.
Then $x|D_a(g)=g^r\,\iota y$ lies in $S$, so merely $g^ex\in S(a)$,
and $\lambda^S((g^ex)/g^{r+e})$ has image $\lambda^E(x/g^r)=\iota y$.
Injectivity of $\iota$ gives $\lambda^S((g^ex)/g^{r+e})=y$. The fibre
is a proposition, so the mere exponent suffices.
\end{proof}

\begin{lemma}[Localisation of the support ideal]\label{lem:J-qc}
For every $a$ and $g\in\OO(a)$, if $\lambda^{\OO}_{a,g}$ is an
isomorphism then so is $\lambda^{J_U}_{a,g}$. In particular $J_U$
satisfies the $\lambda$-criterion on every affine, and is quasi-coherent.
\end{lemma}
\begin{proof}
$J_U\subseteq\OO$ is saturated with exponent $1$: if
$x\in\OO(a)$ and $D_{D_a(g)}(x|D_a(g))\le U$, then
\[
 D_a(gx)=D_a(g)\wedge D_a(x)=D_{D_a(g)}(x|D_a(g))\le U
\]
by~\eqref{eq:supportops} and~\eqref{eq:restrict-support}, so
$gx\in J_U(a)$. Apply Lemma~\ref{lem:saturation}. On an affine $w$,
$\lambda^{\OO}_{w,s}$ is an isomorphism by the chart ring fact, and
Lemma~\ref{lem:lambda-criterion} gives the model of $J_U|w$ with
coefficients $J_U(w)$.
\end{proof}

No presentation of $w\wedge U$, radical ideal, formal Nullstellensatz
or radical localisation is needed. (One can show $J_U(w)=\sqrt{I_w}$
for a presentation $w\wedge U=\bigvee D_w(h_i)$, but the proof does not
use it.)

\section{A subordinate partition of unity}\label{sec:partition}
\begin{proposition}\label{prop:principalcover}
Under the hypothesis of Theorem~\ref{thm:serre}, for every actual
finite affine cover $u_1,\ldots,u_m$ there merely exist
$f_i\in A=\OO(1)$ with
\begin{equation}\label{eq:global-unit}
 \sum_i f_i=1,\qquad D(f_i)\le u_i.
\end{equation}
Consequently the $D(f_i)$ form a finite global principal affine cover,
and the $D_A(f_i)$ cover $\Spec A$.
\end{proposition}
\begin{proof}
Take $J_i=J_{u_i}$. Lemma~\ref{lem:J-sheaf} gives
$J_i|u_i=\OO|u_i$, so form the partition torsor $T$ and its
translation sheaf $K$ from Lemma~\ref{lem:lifts}.

\emph{The translation sheaf is saturated.}
Use the given cover itself as the quasi-coherence atlas. Let
$E=\prod_iJ_i$ and $q:E\to\OO$ the sum, so $K=\ker q$. Fix an affine
member $w=u_j$ and $s\in\OO(w)$. By Lemma~\ref{lem:J-qc} and the chart
ring fact every $\lambda^{J_i}_{w,s}$ is an isomorphism, so
$\lambda^E_{w,s}$ is one, coordinatewise. The kernel is saturated: if
$x\in E(w)$ and $q(x)|D_w(s)=q(x|D_w(s))=0$, then
$\lambda^{\OO}(q(x)/1)=0$, so injectivity of $\lambda^{\OO}_{w,s}$ gives
$s^eq(x)=q(s^ex)=0$ for some $e$, that is $s^ex\in K(w)$.
Lemma~\ref{lem:saturation} makes $\lambda^K_{w,s}$ an isomorphism, and
Lemma~\ref{lem:lambda-criterion} gives a model of $K|u_j$ with
coefficients $K(u_j)$. These models prove quasi-coherence of $K$ on the
original atlas. There is no splitting, no module-kernel calculation and
no atlas refinement in this step.

\emph{One application of the torsor hypothesis.}
A global point of $T$, supplied merely by the hypothesis, is exactly
a tuple $f_i\in J_i(1)$ with $\sum_i f_i=1$. Its memberships give
$D(f_i)\le u_i$. The support laws imply
\[
 1=D(\sum_i f_i)\le\bigvee_iD(f_i)\le1.
\]
Moreover
\[
 D(f_i)=u_i\wedge D(f_i)=D_{u_i}(f_i|u_i)
\]
is a principal open of an affine, hence affine. The same unit equation
in $A$ gives $\bigvee_iD_A(f_i)=1$ in its Zariski lattice.
For the empty cover the unique empty tuple gives $1=0$ in $A$;
$X$ and $\Spec A$ are both empty. All conclusions remain under the
mere existence assertion.
\end{proof}

This simultaneous lift replaces the quotient-supported construction,
affine replacement, its finite induction and the later row-map lift.
The only quasi-coherence argument for its translation sheaf is
saturation (Lemma~\ref{lem:saturation}), used twice.

\section{From that cover to affineness}\label{sec:finish}
We first isolate the structure-ring comparison needed for the final
canonical map. Its hypotheses are the particular cover just constructed.

\begin{lemma}[Ring localisation on a principal affine cover]
\label{lem:principal-ring-localisation}
Suppose $\bigvee_i D(f_i)=1$, with each $D(f_i)$ affine and
$f_i\in A=\OO(1)$. Then for each member of the cover the canonical
ring map
\[
 \lambda_k:A_{f_k}\longrightarrow\OO(D(f_k))
\]
is an isomorphism.
\end{lemma}
\begin{proof}
Fix $k$ and use Theorem~\ref{thm:localisation-descent} for $M=\OO$
and $g=f_k$. Each member $U_i=D(f_i)$ is affine. Each overlap
\[
 U_i\wedge U_j=D_{U_i}(f_j|U_i)
\]
is a principal open of an affine, hence affine too. On either of these
affines $w$, the canonical chart ring fact
\[
 \OO(w)_{f_k|w}\cong\OO(D_w(f_k|w))
\]
is exactly the required component $\lambda$. No associated model
of $\OO$ is used. Pure descent makes the global canonical module
map an isomorphism.

Restriction $A\to\OO(D(f_k))$ is a ring map inverting $f_k$,
so it has a unique ring extension from $A_{f_k}$. Its underlying
module map is the same canonical localisation map, by uniqueness
on original elements. Hence this specified ring extension is a ring
isomorphism.
\end{proof}

Merely choose a finite affine atlas and apply
Proposition~\ref{prop:principalcover}. Its single partition-of-unity
equation already gives both covers: $D(f_i)$ of $X$ and $D_A(f_i)$
of $\Spec A$, where $A=\OO(1)$. Apply the preceding lemma to obtain
\begin{equation}\label{eq:ringlocal}
 A_{f_i}\cong\OO(D(f_i))
\end{equation}
by the canonical ring extensions of restriction.
The canonical geometric map
$\gamma:X\to\Spec A$ has inverse-image map
$D_A(f)\mapsto D(f)$, and its ring map on these opens is precisely
\eqref{eq:ringlocal}. Since $D(f_i)$ is affine, this identifies
the restricted map
\[
 \gamma_i:X|D(f_i)\longrightarrow(\Spec A)|D_A(f_i)
\]
with an isomorphism of affine schemes. Agreement of the canonical
map with the spectrum of the ring map follows from the affine
universal property; it is not an unmentioned computation.

\begin{lemma}[Isomorphisms are local on the target]\label{lem:localiso}
Let $g:Y\to Z$ be a morphism of locally ringed lattices and
$z_i$ a finite cover of $Z$. If each map
$Y|g^*z_i\to Z|z_i$ is an isomorphism, then $g$ is an isomorphism.
\end{lemma}
\begin{proof}
For $y\in L_Y$, let $b_i\le z_i$ be the inverse image of
$y\wedge g^*z_i$ under the lattice isomorphism on the $i$th
restriction. On intersections these inverse lattice maps agree,
because both invert the restricted $g^*$; hence
$b_i\wedge z_j=b_j\wedge z_i$.
Put $b=\bigvee_i b_i$. Distributivity gives $b\wedge z_i=b_i$,
and therefore $g^*b=y$, since the $g^*z_i$ cover $Y$.
Similarly equality of $g^*b$ and $g^*b'$ implies equality after
meeting with every $z_i$, hence $b=b'$. Thus $g^*$ is a lattice
isomorphism.

For every open $z$, its cover by $z\wedge z_i$ and the source
cover by their inverse images give sheaf equalisers. The maps on
all member and overlap rings are isomorphisms by the local
isomorphism hypotheses and restriction. These induce an inverse
on the equaliser ring map. The inverses are natural by uniqueness,
and form the inverse ringed morphism. Isomorphisms preserve
invertibility opens by their universal characterisation, so the
inverse is local. This gives the required geometric isomorphism.
\end{proof}

The lemma applied to $\gamma$ and the target cover given by
\eqref{eq:global-unit} proves Theorem~\ref{thm:serre}.
All remaining witnesses are eliminated into the proposition that
this specified canonical map is an isomorphism. If the cover is
empty, $X$ is the empty lattice space, $A$ is the zero ring, and
the same empty-cover argument gives its affine identification.

Conversely an affine scheme satisfies the hypothesis by
Lemma~\ref{lem:affine-torsor}. Thus this proves an equivalence
between affineness and quasi-coherent first-cohomology vanishing
for the schemes considered here.

\section{Formalisation}\label{sec:review}
\sloppy
Everything above, except the optional module theory of
Section~\ref{sec:modules} and the restriction and locality of
quasi-coherence (replaced by Corollary~\ref{cor:lambda-everywhere}), is
checked in \code{unit/sag/qcqs} (30 October 2026). No new axioms are
introduced; the development rests on the Theory parameters univalence
and replacement. The main statements are:
\begin{itemize}
\item direct implication: \code{SerreAffineness.Serre[L,X,qcqs,hypothesis].affine},
 with hypothesis \code{SerrePartitionCover.On[L,X].Hypothesis} (every
 torsor under every small quasi-coherent module merely has a global
 section) and conclusion \code{AffineCriterion.IsAffine};
\item converse and equivalence: \code{SerreAffineConverse.On[L,X,affine].hypothesis},
 \code{SerreAffineConverse.Criterion[L,X,qcqs]};
\item $\HH^1$ form: \code{SerreCechVanishing.On[L,X].Theorem[qcqs]}
 (\code{affine-to-every}, \code{every-to-some}, \code{some-to-affine});
\item independence of the affine cover:
 \code{SerreCechComparison.On[L,X,M,qc].Covers[U,V]}
 (\code{equivalence}, \code{comparison-add}, \code{comparison-torsor}).
\end{itemize}
The formal proof follows the text, one module per step:
\begin{center}\small
\begin{tabular}{ll}
canonical maps $\lambda_{a,g}$, element form & \code{SerreLambda}, \code{SerreLambdaClearing},\\
 & \code{SerreLambdaElementarySquare}\\
chart ring fact & \code{SerreLambdaChart}\\
Lemma~\ref{lem:lambda-criterion}, $\lambda\Rightarrow$ model & \code{SerreLambdaCriterion}\\
Lemma~\ref{lem:lambda-criterion}, model $\Rightarrow\lambda$ & \code{SerreLambdaModel}\\
finite products & \code{SerreLambdaProducts}\\
Theorem~\ref{thm:localisation-descent} & \code{SerreLocalisationDescent},\\
 & \code{SerreLambdaTop}, \code{SerreLambdaNested}\\
Corollary~\ref{cor:affine-basic-sections} & \code{SerreAffineLocalisation}\\
Corollary~\ref{cor:lambda-everywhere} & \code{SerreLambdaQuasiCoherent}, \code{SerreLambdaUnit}\\
Lemma~\ref{lem:h1affine} & \code{SerreModuleCechGluing}\\
Lemma~\ref{lem:associated-affine-torsor} & \code{SerreAffineTorsor}\\
Lemma~\ref{lem:torsor-chart} & \code{SerreTorsorChart}\\
Corollary~\ref{lem:affine-torsor} & \code{SerreAffineConverse}\\
Proposition~\ref{prop:classification} & \code{SerreCechIsoClass}\\
Corollary~\ref{cor:coverindependence} & \code{SerreCechVanishing}\\
Lemma~\ref{lem:refinement} & \code{SerreCechRefinement}\\
Corollary~\ref{cor:groupcomparison} & \code{SerreCechComparison}\\
Lemma~\ref{lem:lifts} & \code{SerrePartition} (via \code{SerreTorsorLifts})\\
Lemma~\ref{lem:J-sheaf} & \code{SerreSupportIdeal}\\
Lemmas~\ref{lem:saturation}, \ref{lem:J-qc} & \code{SerreSupportIdealSaturation},\\
 & \code{SerreSupportIdealLambda}\\
Proposition~\ref{prop:principalcover} & \code{SerrePartitionQuasiCoherent}, \code{SerrePartitionCover}\\
Lemma~\ref{lem:principal-ring-localisation} & \code{SerreRingLocalisation}\\
Lemma~\ref{lem:localiso} & \code{SerreLocalIso}\\
the restricted canonical map & \code{SerreCanonicalMap}\\
Theorem~\ref{thm:serre}, direct & \code{SerreAffineness}
\end{tabular}
\end{center}

\subsection*{Simplifications made for the formal proof}
\begin{enumerate}
\item Saturation (Lemma~\ref{lem:saturation}) gives the localisation of
 the support ideal and of the partition kernel; no radical description
 of $J_U$ on an affine and no splitting of the kernel are used.
\item In Lemma~\ref{lem:localiso} the relation
 $b_i\wedge z_j=b_j\wedge z_i$ is not needed: surjectivity uses only
 $g^*(\bigvee_ib_i)=y$.
\item The direction model $\Rightarrow\lambda$ of
 Lemma~\ref{lem:lambda-criterion} is proved in the element form, on
 every principal open of the chart, so no restricted model and no
 associated restriction comparison are needed.
\item Corollary~\ref{cor:lambda-everywhere} replaces the restriction of
 quasi-coherence to affine opens.
\item Affine torsor triviality goes through affine reconstruction and
 Lemma~\ref{lem:torsor-chart} to the associated case, not through a
 second degree-one argument.
\item The $\HH^1$ form of the theorem is stated by vanishing
 (Corollary~\ref{cor:coverindependence}); the comparison of the groups
 (Corollary~\ref{cor:groupcomparison}) goes through torsors and the
 concatenation of the two covers, with additivity from restriction
 instead of a contracted sum of torsors.
\end{enumerate}

\subsection*{Validation}
The full server check has 390 modules and no failures (16.4\,s). The
independent checker, run on the whole development with opening
instrumentation, checks 16{,}047 declarations (35.1\,s). Every one of
the 6{,}797 declarations of the 141 Serre modules has exactly one audit
entry stating which definitions it may open, and every one of the
10{,}549 entries of the combined audit reports that nothing else was
opened.

\begin{thebibliography}{9}
\bibitem{StacksSerre}
The Stacks Project, Lemma 30.3.1, Serre's criterion for affineness,
\href{https://stacks.math.columbia.edu/tag/01XF}{Tag 01XF}.
Its closed-point proof is a comparison reference, not used here.
\bibitem{StacksTorsors}
The Stacks Project, Lemma 21.10.4, Cech cohomology and torsors,
\href{https://stacks.math.columbia.edu/tag/0A6G}{Tag 0A6G}.
The explicit finite-cover construction in this note gives the
needed degree-one comparison without derived cohomology.
\end{thebibliography}
\end{document}
